\section{Conclusion}

Multi-agent systems increasingly rely on protocols such as MCP and A2A.
However, current approaches tend to expose execution-level details directly into the agent's planning context. 
This forces agents to reason simultaneously about task semantics and protocol mechanics, increasing token consumption, latency, and failure rates.

In this work, we formalize this as a problem of poorly staged information exposure and address it with a unified capability gateway that abstracts over both tools and peer agents. 
The gateway keeps planning focused on semantic capability information and delays execution-specific details until they are actually needed.
Our empirical evaluation showed that the gateway architecture consistently improves token efficiency, latency, and task success. 
The retrieval-based gateway variant becomes especially useful as the environment grows, preventing irrelevant information from entering the context window.

Taken together, these results suggest that future agent systems should be designed around staged information exposure rather than upfront capability dumping.
The gateway architecture reinforces this principle while additionally providing a central point for configuration and governance across protocols.
As multi-agent ecosystems continue to grow, this separation should become increasingly important for building reliable and scalable systems.